\documentclass[11pt]{article}
\usepackage[margin=1in]{geometry}
\usepackage{amsmath,amssymb}
\usepackage{graphicx}
\usepackage{booktabs}
\usepackage{hyperref}
\usepackage{natbib}

\title{The SUMO E3 Ligase Ema2 Promotes Genome-Wide lncRNA Transcription from the Somatic Macronucleus to Enable scnRNA-Directed DNA Elimination in \textit{Tetrahymena}}
\author{Anonymous Authors}
\date{}

\begin{document}
\maketitle

\begin{abstract}
Small RNA-directed chromatin silencing requires nascent long non-coding RNA (lncRNA) transcription from target loci, but the molecular mechanisms promoting this transcription remain poorly understood. In \textit{Tetrahymena thermophila}, scan RNAs (scnRNAs) produced from the germline micronucleus (MIC) guide the elimination of approximately 12,000 internal eliminated sequences (IESs) from the somatic macronucleus (MAC) during conjugation. Here we identify Ema2 (TTHERM\_00113330) as a conjugation-specific SUMO E3 ligase that is required for genome-wide lncRNA transcription from the parental MAC. Ema2 contains a conserved SP-RING domain, interacts directly with the SUMO E2 enzyme Ubc9, and promotes SUMOylation of the transcription regulator Spt6. Loss of Ema2 blocks target-directed scnRNA degradation (TDSD), impairs heterochromatin formation in the new MAC, and partially disrupts DNA elimination, producing inviable sexual progeny. Ema2 is required for the chromatin association of both Spt6 and RNA polymerase II (Rpb3) during mid-conjugation, providing a mechanistic link between SUMOylation and lncRNA transcription. However, SUMOylation-defective Spt6 mutants do not phenocopy the lncRNA defect, indicating that Ema2 has additional critical targets beyond Spt6. Our findings reveal that SUMOylation actively promotes the lncRNA transcription that enables small RNA-mediated genome surveillance.
\end{abstract}

\section{Introduction}

Programmed DNA elimination is a widespread epigenetic phenomenon in which specific sequences are removed from the somatic genome during development \citep{yao1984programmed}. In the ciliate \textit{Tetrahymena thermophila}, this process eliminates approximately 97~Mb of DNA ($\sim$12,000 IESs) from the newly forming macronucleus (MAC) after sexual conjugation \citep{mochizuki2002analysis}. The specificity of elimination is guided by scan RNAs (scnRNAs), a class of Piwi-interacting small RNAs produced from the germline micronucleus (MIC) during meiosis.

The scnRNA pathway involves a critical scanning step: scnRNAs loaded onto the Piwi protein Twi1 are compared against lncRNAs transcribed from the parental MAC \citep{aronica2008studying}. ScnRNAs complementary to MAC sequences are selectively degraded through target-directed scnRNA degradation (TDSD), leaving only IES-matching scnRNAs to guide heterochromatin formation and DNA elimination in the new MAC \citep{noto2015tetrahymena}. This creates a paradox: the somatic genome destined for reorganization must be actively transcribed to produce the lncRNAs that mediate genome scanning.

The molecular mechanism that actively promotes this lncRNA transcription from the parental MAC during the critical mid-conjugation window was unknown. SUMOylation---the covalent attachment of SUMO proteins to target lysines, catalyzed by E1, E2, and E3 enzymes \citep{johnson2004sumo}---regulates many chromatin processes, but its role in lncRNA transcription for genome surveillance had not been established. The transcription elongation factor Spt6 \citep{endoh2009spt6} is essential for transcription, but how its activity is modulated specifically for lncRNA production during conjugation was unclear.

\section{Results}

\subsection{Ema2 Is a Conjugation-Specific MAC-Localized Protein}

EMA2 mRNA was undetectable in vegetative and starved cells but was highly expressed during conjugation (3--12 hours post-mixing, hpm). Ema2-HA protein first appeared in the parental MAC at 3 hpm, then switched to the new MAC at 8 hpm before fading by 14 hpm. This localization switch from parental to new MAC is reminiscent of the Piwi protein Twi1 and PRC2 components.

\subsection{Ema2 Is Required for DNA Elimination and Progeny Viability}

FISH analysis with Tlr1 transposable element probes at 36 hpm showed that EMA2 knockout exconjugants retained Tlr1 signal in new MACs, indicating partially blocked DNA elimination. The Tlr1 signal was weaker than in TWI1 knockout controls, suggesting partial rather than complete elimination failure. EMA2 KO crosses produced no viable sexual progeny.

\subsection{Ema2 Is Required for Target-Directed scnRNA Degradation}

Northern blot analysis with MDS-complementary probes showed that in wild-type cells, MDS-matching scnRNAs were progressively degraded between 3 and 6 hpm (TDSD). In EMA2 KO cells, these scnRNAs persisted, demonstrating that TDSD was blocked. H3K27me3 heterochromatin marks in the new MAC were reduced or absent in the KO, consistent with failure of downstream chromatin modification.

\subsection{Ema2 Functions as a SUMO E3 Ligase}

Sequence analysis revealed a conserved SP-RING domain in Ema2, with zinc-coordinating residues aligned with known SUMO E3 ligases (MMS21, SIZ1, Su(var)2-10, PIAS1). In vitro pulldown assays demonstrated direct interaction between recombinant GST-Ema2 and His-Ubc9 (SUMO E2 enzyme). Addition of Ema2 to E1+E2 reactions enhanced SUMOylation activity, confirming E3 ligase function.

\subsection{Ema2 Promotes SUMOylation of Spt6}

His-Smt3 pulldown followed by mass spectrometry from conjugating cells at 6 hpm identified Spt6 as a protein with reduced SUMOylation in EMA2 KO compared to wild-type. Western blots confirmed slower-migrating Spt6-HA species appearing during conjugation (4.5--6 hpm), consistent with SUMOylated forms.

\subsection{Ema2 Is Required for lncRNA Accumulation}

RT-PCR using primers spanning MDS-IES junctions detected lncRNAs from the parental MAC in wild-type cells at 6 hpm but not in EMA2 KO cells. Immunofluorescence with the J2 anti-dsRNA antibody confirmed strong dsRNA signal in wild-type parental MACs that was absent or strongly reduced in the KO. Genomic DNA controls confirmed that target loci were intact in KO strains.

\subsection{Ema2 Promotes Chromatin Association of Spt6 and RNA Polymerase II}

Immunofluorescence showed that both Spt6-HA and RNAPII subunit Rpb3 had reduced chromatin association in EMA2 KO compared to wild-type crosses during mid-conjugation. The majority of Rpb3 signal was lost from chromatin, indicating impaired transcription engagement.

\subsection{SUMOylation-Defective Spt6 Does Not Phenocopy EMA2 KO}

A critical control experiment showed that SUMOylation-defective Spt6 mutants maintained normal lncRNA production and DNA elimination. This indicates that Spt6 SUMOylation alone is insufficient to explain Ema2's role, and that additional SUMOylation targets beyond Spt6 are critical for lncRNA transcription.

\section{Discussion}

Our findings establish that the SUMO E3 ligase Ema2 actively promotes genome-wide lncRNA transcription from the parental MAC, connecting SUMOylation to small RNA-mediated genome surveillance. The pathway proceeds through promotion of Spt6 and RNAPII chromatin association, but the failure of SUMOylation-defective Spt6 to phenocopy EMA2 KO indicates that multiple Ema2 targets contribute to lncRNA transcription.

\textit{Tetrahymena} provides a uniquely tractable system for studying this process because the source (MIC) and target (MAC) of small RNAs reside in different nuclei, enabling clean dissection of lncRNA transcription mechanisms that would be confounded in organisms where source and target loci overlap.

The conjugation-specific expression and MAC-to-new-MAC localization switch of Ema2 ensures that SUMOylation-dependent lncRNA transcription is restricted to the appropriate developmental window. This temporal control adds to the growing understanding of how developmental timing coordinates the multiple steps of programmed DNA elimination.

\section{Methods}

Somatic knockout strains were generated by disrupting all MAC copies of EMA2. Ema2-HA was tagged at the endogenous locus for localization studies. DNA elimination was assayed by FISH with Tlr1 probes at 36 hpm. TDSD was monitored by northern blot with MDS-complementary probes. In vivo SUMOylation was analyzed by Ni-NTA pulldown of His-tagged Smt3 conjugates followed by mass spectrometry. Chromatin association was assessed by immunofluorescence with anti-Spt6-HA and anti-Rpb3 antibodies. lncRNA accumulation was measured by RT-PCR with MDS-IES junction primers and J2 antibody immunofluorescence.

\bibliographystyle{plainnat}
\bibliography{references}

\end{document}
